\section{Well-Bracketed OGS}
\label{s:wb-ogs}

In the operational game semantics we have studied so far, Opponent is allowed to 
answer over continuation names in an arbitrary order, and is 
even allowed `to forget' to answer 
to some of these names.
This behavior of Opponent 
corresponds to that of a context that has 
access to control operators, e.g. $\mathtt{call/cc}$, with an affine usage of continuations.


\subsection{Well-Bracketed LTS}


\begin{figure*}
\[\begin{array}{lllll}
 (P\tau) & \conf{M,p,\stt,\gamma,\phi} & \arrWB{\ \tau \ } & 
  \conf{N,p,\stt,\gamma,\phi} & \text{ when } M \red N\\
%
 (PA) & \conf{V,p,\stt,\gamma,\phi} & \arrWB{\ansP{p}{x}} & 
  \conf{\stt,\gamma \cdot \pmap{x}{V},\phi \uplus \{x\}} \\
%
 (PQ) & \conf{E[xV],p,\stt,\gamma,\phi} & \arrWB{\questPV{x}{y}{q}} & 
  \multicolumn{2}{c}
    {\conf{q::\stt,\gamma \cdot \pmap{y}{V}\cdot \pmap{q}{(E,p)},\phi \uplus \{y,q\}}} \\
%
 (OA) & \conf{p::\stt,\gamma\cdot [p \mapsto (E,q)],\phi} & \arrWB{\ansO{p}{x}} & 
  \conf{E[x],q,\stt,\gamma,\phi \uplus \{x\}} \\
%
 (OQ) & \conf{\stt,\gamma,\phi} & \arrWB{\questOV{x}{y}{p}} & 
  \conf{V y,p,\stt,\gamma,\phi \uplus \{y,p\}} & 
  \text{ when } \gamma(x) = V\\
%
 (IOQ) & \initconf{\phi}{M} & \arrWB{\questOinit{p}} &
     \conf{M,p,[],\emptymap,\phi \uplus \{p\}}
\end{array}\]
\caption{Well-Bracketed OGS LTS}
\label{fig:stacked-lts}
\end{figure*}

To capture the absence of such control operators in the environment,
we enforce a `well-bracketed' behavior on the sequence of answers.
To do so, in Figure~\ref{fig:stacked-lts} we introduce the  \emph{Well-Bracketed OGS} (\WBOGS).  
It refines \AOGS by using \emph{stacked configurations} $\FFb$.
They use a \emph{stack} of continuation names $\stt$, that keeps track
of the order in which the Opponent answers must occur to be well-bracketed.
Such stacks
are represented as lists of Player continuation names.
We write $p:: \stt$ for the stack with $p$ pushed on top of $\stt$, 
and $\emptystack$ for the empty stack.
In \WBOGS, Player questions push their fresh continuation names on top of the stack, and
Opponent answers pop the stack to obtain the continuation name they are allowed to answer to.

Extending the notion of valid configurations from \AOGS,
we require the stack to be formed by distinct Player continuation names.
Since names in stacks are all distinct, removing the evaluation context 
from $\gamma$ after an Opponent answer is not necessary anymore, since the 
associated continuation name will never be allowed to be used again. 
But for uniformity \wrt the other LTS, we choose to keep it.
For a stacked configuration $\FFb$, we write $\es \FFb$ for the corresponding \AOGS
configuration, where the stack has been removed.

Notice that stacked configuration $\FFb$
only stores the stack of Player continuations as an independent component,
in order to restrict Opponent behavior.
By shuffling the continuation structure of $\gamma$ with
the stack $\sigma$, we could reconstruct
a stack for the whole interaction between Player and Opponent.

\subsection{Well-Bracketed Traces}

The goal is now to characterize the set of traces generated by a stacked configuration
$\FFb$ (on the well-bracketed LTS) as the subset of the traces generated by $\es \FFb$ 
(on the alternating LTS) that satisfies the following condition of \emph{well-bracketing}.

Considering a trace $\tr_1 \ \act_1  \ \tr_2 \ \act_2 \ \tr_3$,
with $\act_2$ an answer of the shape $\ansO{p}{x}$ or $\ansP{p}{x}$
and $\act_1$ the action that introduces the continuation name $p$,
then we say that $\act_1$ is \emph{answered} by $\act_2$.
Since only questions can introduce continuation names, 
$\act_1$ is indeed a question, and it must be of the opposite polarity
as $\act_2$.
%
Well-bracketing corresponds to the fact that questions must be answered in the last-in
first-out order they previously appear in a trace.
We formalize this idea using a pushdown system, defined as an LTS whose configurations are 
stacks of distinct continuation names with transition relation defined as:
\begin{mathpar}
  (\text{PQ}) \
    {\stt \pshdtrans{\questP{x}{y,p}} p::\stt }
  \and
  (\textsc{OQ}) \
    {\stt \pshdtrans{\questO{x}{y,p}} \stt }
  \and
  (\textsc{PA}) \
    {p::\stt \pshdtrans{\ansP{p}{x}} \stt }
  \and
  (\textsc{OA}) \
    {p::\stt \pshdtrans{\ansO{p}{x}} \stt }
\end{mathpar}

\begin{definition}
Taking a stack $\stt$,
a trace $\tr$ is said to be \emph{well-bracketed} \wrt $\stt$ when
there exists another stack $\stt'$ such that
$\stt \pshdTrans{\tr} \stt'$.
\end{definition}
%
For example $\ansO{c}{w} \questP{x}{y,c'} \ansO{c'}{z} \ansP{c''}{w}$
is well-bracketed \wrt $c::[]$.
Notice that we do not keep track of the Opponent continuation names in this definition.
Indeed, Proponent behaviour is automatically well-bracketed,
since running terms are taken only in the pure $\lambda$-calculus.

\begin{lemma}
\label{l:WBs_lts}
Taking $\FFb$ a configuration whose stack is $\stt$, then
$\FFb \ArrWB{\tr} \FFb'$ iff 
$\es \FFb \ArrA{\tr} \es \FFb'$ and $\tr$ is well-bracketed \wrt $\stt$.
\end{lemma} 

\begin{proof}
  Suppose that $\FFb \ArrWB{\tr} \FFb'$,
  then writing $\stt'$ for the stack of $\FFb'$,
  we prove by a direct induction on $\tr$ that $\stt \pshdTrans{\tr} \stt'$.

  In the reverse direction, suppose that 
  $\es \FFb \ArrA{\tr} \FF'$ and $\stt \pshdTrans{\tr} \stt'$,
  we consider the stacked configuration $\FFb'$ built 
  from $\FF'$ using $\stt'$.
  Then we prove by a direct induction on $\tr$ that
  $\FFb \ArrWB{\tr} \FFb'$.
\end{proof}

\subsection{Relationship between Alternating and Well-Bracketing OGS}
\label{subsec:relate-alt-wb-lts}

In Section~\ref{subsec:relate-conc-alt-lts},
we have related the denotation of $\lambda$-terms
given by the \COGS and \AOGS LTS.
We now establish a similar connection
between \AOGS and \WBOGS.
This is done by adapting the proof of Lemma~\ref{l:OGS_COGS}.
To do so, we first adapt the definition of
triggerable singleton subconfigurations.

\begin{definition}
A singleton \AOGS configuration $\FF$ is 
a \emph{triggerable} subconfiguration of a stacked configuration $\FFb$ when:
\begin{itemize}
  \item either $\FFb$ is active of the shape
  $\conf{M,p,\stt,\gamma,\phi}$,
  and $\FF = \conf{M,p,\emptymap,\phi \backslash{\dom{\gamma}}}$;
  \item or $\FFb$ is passive of the shape
  $\conf{\stt,\gamma,\phi}$,
  and $\FF$ is:
  \begin{itemize}
    \item either of the shape
    $\conf{\pmap{x}{V},\phi \backslash{\dom{\gamma}} \cup\{x\}}$
    with $\pmap{x}{V} \subseteq \gamma$;
    \item or of the shape $\conf{\pmap{p}{(E,q)},\phi \backslash{\dom{\gamma}} \cup\{p\}}$
    with $\pmap{p}{(E,q)} \subseteq \gamma$,
    when $p$ is the top element of $\stt$.
  \end{itemize}
\end{itemize}
\end{definition}
\noindent
So
$\conf{\pmap{p}{(E,q)},\phi}$ is not triggerable 
when $p$ is not the top element of the stack,
since it cannot be directly triggered by an OA action of the \WBOGS LTS
(while it can be by the \AOGS LTS).

We recall that we only consider well-formed relations, 
as introduced in Definition~\ref{d:compa}.
Therefore, whenever two configurations $\FF,\GG$ appear as
a pair in a relation (i.e., a behavioural relation) 
it is implicitly assumed that they are support-equivalent, written
$\okC{\FF}{\GG}$.
For relations on stacked configurations, $\okC{\FFb}{\GGb}$ 
also enforces that $\FFb,\GGb$ have the same stack.

\begin{lemma}
\label{l:AOGS_WBOGS}
If $\FFb, \GGb$ are singleton stacked configurations and 
$\FFb \wbWB \GGb$, then also $\es\FFb \wbA \es\GGb$. 
\end{lemma}

\begin{proof}
  Let $\R$ be the relation on \AOGS singleton configurations
  defined as $\FF \RR \GG$ when
   $\okC \FF\GG$ and there exist stacked configurations $\FFb, \GGb$
    with $\FFb \wbWB \GGb$ such that 
    $\FF$ and $\GG$ are triggerable singleton configurations
    respectively of $\FFb$ and $\GGb$.

  First notice that if $\FFb, \GGb$ are stacked singleton configurations and 
  $\FFb \wbWB \GGb$, we directly get that 
  $\es\FFb \RR \es\GGb$ since $\es\FFb$ and $\es\GGb$ are 
  triggerable subconfigurations of $\FFb$ and $\GGb$ respectively.
  
  We now show that $\R$ is a bisimulation up-to composition.
  We take $\FF \RR \GG$, such that there exists $\FFb, \GGb$
  as stated by the definition of $\R$.
  Suppose 
  \begin{equation}
  \label{e:arrA-left}
   \FF \arrA{\act} \FF_1 
  \end{equation} 
  in \AOGS.  
  First, we suppose that $\FF,\GG$ are both passive, 
  so that the transition is either OA or OQ.  
  We assume it is OA, since it is the hardest case, 
  Opponent answers being constrained in \WBOGS.
  We thus have, for some $p,V,x,y$ 
  \[\begin{array}{rcl}
   \FF &=&  \conf{ \gamma_1, \phi} \mbox{ with }  \gamma_1 = \pmap{p}{(E,q)} \\
   \act & = & \ansO{p}{x} \\
   \FF_1 & = & \conf{E[x],q, \emptymap,\phi \uplus \{x\}}
  \end{array}\]
  Moreover, we have $\GG = \conf{\delta_1, \phi}$ 
  with $\delta_1 = \pmap{p}{(\Ep,q)}$ for some evaluation context $\Ep$.
  Since $\FF,\GG$ are triggerable respectively from $\FFb,\GGb$,
  we get that both $\FFb,\GGb$ are passive
  and that the top element of the stack of $\FFb$ is $p$.
  From $\FFb \wbWB \GGb$, we get that $\okC{\FFb}{\GGb}$,
  so there exists a stack $\stt$ and two environments $\gamma, \delta$ sharing the same domain such that:  
  \[\begin{array}{rcl}
  \FFb & = & \conf{p::\stt,\gamma_1 \cdot \gamma,\phi \uplus \dom{\gamma}} \\ 
  \GGb & = & \conf{p::\stt,\delta_1 \cdot \delta,\phi \uplus \dom{\delta}} 
  \end{array}\]
  We have 
  $\FFb \arrWB{\act} \conf{E[x],q,\sigma, \gamma, \phi\uplus \dom{\gamma} \uplus \{x\}} 
  \defi \FFb'$. 
  Since $\FFb \wbWB \GGb$, 
  \[\GGb  \arrWB{\act} \conf{E'[x],q,\sigma, \delta, \phi \uplus \dom{\delta} \uplus \{x\}}  
  \defi \GGb'\]
  and $\FFb' \wbWB \GGb'$.
  %
  Let $\GG_1 = \conf{\Ep[x],q,\delta_1,\phi}$,
  we also have 
  $\GG \arrA \act \GG_1$.
  This is sufficient to match \reff{e:arrA-left}, as 
  $\FF_1 \uptoComp \R  \GG_1$. Indeed we have:
  \begin{itemize}  
  \item $\conf{\gamma_1,\phi\uplus \{y,p\}} \R  \conf{\delta_1,\phi\uplus \{y,p\}}$
  since $\FF \RR \GG$;
  \item  $\conf{E[x],q, \phi\uplus \{x\}} \R  \conf{\Ep[x],q, \phi\uplus \{x\}}$
  since they are triggerable subconfigurations respectively of $\FFb'$ and $\GGb'$.
  \end{itemize}
\noindent
  Now suppose that $\FF,\GG$ are active,
  then the proof is similar to the one for Lemma~\ref{l:OGS_COGS}. \qedhere
  % so there are 3 possibilities of transitions, corresponding to rules 
  % P$\tau$, PA, PQ. The case of 
  % P$\tau$ is straightforward since bisimilarity is preserved by internal
  % actions.  We only look at PQ, as PA is simpler. 
  % Thus we have 
  % $\FF = \conf{E[x V],p,\emptymap,\phi}$, for some $E,x,V$ and:
  % \begin{equation}
  % \label{e:KpqH}
  %  \FF  \arrA{\questPV{x}{y}{q}} 
  %       \conf{\pmap{y}{V} \cdot \pmap{q}{(E,p)}, \phi \uplus \{y,q\}} 
  %       \defi \FF_1
  % \end{equation} 
  % Since $\GG$ shares the same continuation structure as $\FF$,
  % we have $\GG = \conf{M,p,\emptymap,\phi}$ for some term $M$.
  % Since $\FF,\GG$ are triggerable respectively from $\FFb,\GGb$,
  % there exists a stack $\stt$ and two environments $\gamma,\delta$ sharing the same domain such that:  
  % \[\begin{array}{rcl}
  %   \FFb & = & \conf{E[x V], p,\stt, \gamma, \phi\uplus\dom{\gamma}} \\ 
  %   \GGb & = & \conf{M, p,\stt, \delta, \phi\uplus\dom{\delta}} 
  % \end{array}\]
  % We also have 
  % \[
  % \FFb \arrWB{\questPV{x}{y}{q}} 
  %  \conf{q::\stt,\pmap{y}{V} \cdot \pmap{q}{(E,p)} \cdot \gamma, \phi  \uplus \{y,q\}}
  %  \defi \FFb'
  % \]
  % As $\FFb \wbWB \GGb$, there exist $\Ep,W$ such that 
  % \[
  % \GGb \Longrightarrow 
  %  \conf{\Ep[x W],p,\stt,\delta,\phi}
  % \arrWB{\questPV{x}{y}{q}} 
  % \conf{q::\stt,\pmap{y}{W} \cdot \pmap{q}{(\Ep,p)} \cdot \delta, \phi \uplus \{y,q\}} 
  % \defi \GGb'
  % \]
  % with
  % $\FFb' \wbWB \GGb'$. 
  % Defining $\GG_1$ as
  % $\conf{\pmap{y}{W} \cdot \pmap{q}{(\Ep,p)}, \phi \uplus \{y,q\}}$,
  % we also have $\GG \ArrA{\questPV{x}{y}{q}} \GG_1$. 
  % We deduce that both 
  % $\conf{ \pmap{y}{V},\phi \uplus \{y\}} 
  %  \RR
  %  \conf{ \pmap{y}{W},\phi \uplus\{y\}}$ 
  % and 
  % $\conf{ \pmap{q}{(E,p)},\phi \uplus \{q\}}
  %  \RR
  %  \conf{ \pmap{q}{(\Ep,p)},\phi \uplus \{q\}}$ hold
  %  since all these singleton configurations
  %  are triggerable from either $\FFb'$ or $\GGb'$,
  %  and that $\FFb' \wbWB \GGb'$.
  % Hence $\FF_1 \uptoComp \R  \GG_1$.
  
  % In summary, as an answer to the challenge 
  % \reff{e:KpqH}, we have found  $\GG_1$ such that 
  %  $\GG \ArrA{\questPV{x}{y}{q}} \GG_1$ and  
  %  $\FF_1 \uptoComp \R  \GG_1$; this closes the proof.
\end{proof} 

\reviewComment{Lemma 10.4: this statement would be better appreciated with an
example. It seems to me that the point is that the whole "syntax tree"
(enf) is already explored following the wb discipline, because the
P-views are well-bracketed. Is this what is going on?}
\reviewCommentDS{I am not sure i'd like to add an example. 
I would simply add a sentence or two here, confirming the explanation given by the reviewer
}
\reviewCommentGJ{I agree with the beginning of the explanation of the reviewer, that the whole "syntax tree"
(enf) is already explored following the wb discipline. But I'm still trying to understand why the fact that P-views are well-bracketed
is important here. I think it comes from the fact that in the proof, we always work with singleton configurations (thanks to the up-to composition technique), 
which indeed corresponds to only considering the different P-views of a trace.}

\begin{remark}
In the proof of Lemma~\ref{l:AOGS_WBOGS}, we crucially use the fact that the
interaction on singleton configurations is always well-bracketed, as
the stack can have at most one element. Following Remark~\ref{rm:up-to-comp-gs}, this corresponds
to the fact that in game semantics, a P-view is always well-bracketed for Opponent. 
\end{remark}

\subsection{Complete Trace Equivalence}

To relate trace denotation to contextual equivalence, it is important to 
be able to represent terminating interactions between a term and a context.
To do so, 
we consider the notion of \emph{complete trace}, 
that has been introduced in game semantics~\cite{abramsky1997linearity} to  
obtain full-abstraction results for languages with mutable 
store. A trace is said to be complete when all its questions
are answered.
The notion of complete trace is normally considered only for well-bracketed traces,
we generalize it here to more general traces,
and with configurations having already Proponent continuation names.

\begin{definition}
\label{def:ctrace}
 A trace $\tr$ 
 is \emph{complete for an \AOGS (resp. \COGS) configuration $\FF$} if
 $\FF \ArrA{\tr}$ (resp. $\FF \ArrC{\tr}$) 
 and the subject names of the unjustified answers in $\tr$
 are exactly the Player continuation names of $\FF$.
 This definition is extended to a stacked configuration $\FFb$
 with $\stt$ its stack,
 by asking that $\FFb \ArrWB{\tr}$ and $\stt \pshdTrans{\tr} \emptystack$.
 Then we say that a configuration is \emph{terminating} 
 if it has at least one complete trace.
\end{definition}
\noindent
So a trace $\tr$ is complete for $\FF$ when it has no unanswered questions,
and furthermore $\tr$ has also answers to all the pending questions of $\FF$.


\begin{definition}
  \label{d:TEwbcN}
  We write $\FFb_1 \TEwbc \FFb_2$ if
  $\okC{\FFb_1}{\FFb_2}$
  and $\FFb_1,\FFb_2$ have the same set of complete traces in \WBOGS. 
%
  Similarly, we write $\FF_1 \TEc \FF_2$ (resp.\ $\FF_1 \TEpc \FF_2$)
  if $\okC{\FF_1}{\FF_2}$ 
  and $\FF_1,\FF_2$ have the same 
  sets of complete traces in \AOGS (resp.\ \COGS).
\end{definition} 

We now characterize configurations that are reached
by complete traces.

\begin{definition}
  \label{d:sp}
  An \AOGS or stacked configuration
  is \emph{strongly passive} if it is passive and its 
  continuation structure is empty.
  For \COGS, we also ask that its running term be empty.
\end{definition}
\noindent
This means, for stacked configurations,
that the stack is empty, and the configuration is of form
\[   
    \conf{\emptystack, 
      \pmap{x_1}{V_1}\cdot \ldots \cdot  \pmap{x_n}{V_n},\phi} \hskip .8cm n\geq 0 \]
so that indeed we can consider the interaction terminated, since Opponent
cannot answer anymore, even if it can still interrogate one of the $x_i$.

Traces that correspond to reduction to strongly passive configurations 
are complete traces, a result that holds for \WBOGS, \AOGS and \COGS.
  
\begin{lemma}
\label{l:WBOGS_complete-strongly-passive}
  Taking $\FFb,\FFb'$ two stacked configurations 
  and $\tr$ a trace such that $\FFb \ArrWB \tr \FFb'$.  
  Then $\tr$ is complete for $\FFb$
  iff $\FFb'$ is strongly passive.
\end{lemma}

\begin{proof}
  Writing $\stt,\stt'$ for the stack respectively of $\FFb,\FFb'$,
  we get from Lemma~\ref{l:WBs_lts} that
  $\stt \pshdTrans{\tr} \stt'$.
  By definition, $\tr$ is complete for $\FFb$ iff $\stt'$ is empty,
  and $\FFb'$ is strongly passive iff $\stt'$ is empty.
\end{proof}

\begin{lemma}
  \label{l:OGS_complete-strongly-passive}
    Taking $\FF,\FF'$ two \AOGS (resp. \COGS) configurations
    and $\tr$ a trace  s.t. $\FF \ArrA \tr \FF'$
    (resp. $\FF \ArrC \tr \FF'$).  
    Then $\tr$ is complete for $\FF$
    iff $\FF'$ is strongly passive. 
\end{lemma} 

\begin{proof}
  Suppose that $\tr$ is complete for $\FF$.
  All Player questions of $\tr$ having been answered during the interaction, 
  the set of Player continuation names of $\FF'$
  is included in the one of $\FF$.
  Moreover, $\tr$ contains answers for all Player continuation names of $\FF$
  so that no Player continuation name of $\FF'$ is a Player continuation name of $\FF$.
  Therefore, its set of continuation names is empty,
  so that $\FF'$ is strongly passive.

  Reciprocally, if $\FF'$ is passive, then its set of continuation names is empty.
  Thus, $\tr$ cannot have pending questions and must
  have answered all the questions of $\FF$.
  So $\tr$ is complete for $\FF$.
\end{proof}
\noindent
Complete well-bracketed trace equivalence corresponds to contextual equivalence for a 
$\lambda$-calculus with store, as proved in~\cite{abramsky1997linearity} for 
a compositional game model of Idealized Algol,
and in~\cite{laird2007} for an operational game model of RefML.
\cutLMCSsecondRound{
This last result can be adapted to the $\lambdarho$-calculus, 
an untyped call-by-value language with store. Its definition is given in 
Appendix~\ref{a:LambdaRho}, with the definition of 
its contextual equivalence $\ctxeqlr$.
Being a language where the exchange of locations is forbidden,
the proof of full abstraction of the \WBOGS semantics
for $\lambdarho$-calculus can be obtained by a direct adaptation of the result proven in~\cite{JM21}
for the `HOS' language.

\begin{theorem}
\label{t:rho}
Let $M_1,M_2 \in \LasV$
 with free variables in $\phi$.
We have $\initconf{\phi}{M} \TEwbc \initconf{\phi}{N}$
 iff  $M \ctxeqlr N$.
\end{theorem}}
